\scanpage{76}
est naturellement pondérée, c'est $X'$. Si la structure triangulaire $X$ est \inserted{``triangulation} stricte'' (pas de monogone, ni d'arête pliée) alors le revêtement $X'$ est ramifié au plus qu'en les sommets -- et de façon précise \deleted{\unsure{cette}} il y a ramification en un sommet \inserted{de $X$}, sss celui-ci est d'ordre impair. Quand tous les ordres des sommets sont pairs, \unsure{en} plus du fait que toutes les faces sont des triangles et qu'il n'y a pas d'arêtes pliées (\deleted{toutes}\inserted{conditions} nécessaires pour que la structure \inserted{de carte} triangulaire sur $X$ soit pondérable), alors le revêtement envisagé de $X$ est non ramifié. Les sections correspondent $1$-$1$ aux pondérations de la structure triangulaire sur $X$, sauf erreur.

On peut décrire encore un peu plus la situation, dans \deleted{\unsure{le}} le cas orienté qui nous intéresse ici, en disant qu'une condition nécessaire de pondérabilité est l'existence d'un ``coloriage'' en deux ``couleurs'' $+$ et $-$, de façon que deux faces adjacentes soient de couleurs différentes. Les coloriages correspondent aux sections du fibré \inserted{$X'/\unsure{\mathcal G_3|\mathfrak S_3}^+$} associé au fibré de groupe $\unsure{\mathcal G_3|\mathfrak S_3}$, par
\[
\unsure{\mathcal G_3|\mathfrak S_3}\xrightarrow{\mathrm{sg}}\pm1,
\]
\unsure{donc aux} \unsure{restrictions}
\margin{75}
